%% LyX 2.4.4 created this file.  For more info, see https://www.lyx.org/.
%% Do not edit unless you really know what you are doing.
\documentclass[english]{article}
\usepackage[T1]{fontenc}
\usepackage[latin9]{inputenc}
\usepackage{units}
\usepackage{enumitem}
\usepackage{amsmath}
\usepackage{amsthm}
\usepackage{cancel}
\usepackage{geometry}
\geometry{verbose,tmargin=1in,bmargin=1in,lmargin=1in,rmargin=1in}

\makeatletter
%%%%%%%%%%%%%%%%%%%%%%%%%%%%%% Textclass specific LaTeX commands.
\numberwithin{equation}{section}
\numberwithin{figure}{section}
\newlength{\lyxlabelwidth}      % auxiliary length 

%%%%%%%%%%%%%%%%%%%%%%%%%%%%%% User specified LaTeX commands.
\usepackage{fancyhdr}
\usepackage{lastpage}
\pagestyle{fancy}
\usepackage{enumitem}
\usepackage{ifthen}
%sets math fonts to be more readable
\everymath=\expandafter{\the\everymath\displaystyle}
%\DeclareMathSizes{10}{10}{10}{10}
\usepackage{multicol}

\usepackage{tikz}
\usetikzlibrary{plotmarks}
\usepackage{pgfplots}
\pgfplotsset{compat=newest}

\usepackage{calc}
\usepackage[nomessages]{fp}

\usepackage{advdate}    % Advancing/saving dates
\usepackage{datetime}   % Dates formatting
\usepackage{datenumber} % Counters for dates
% Added by lyx2lyx
\setlength{\parskip}{\smallskipamount}
\setlength{\parindent}{0pt}

\makeatother

\usepackage{babel}
\begin{document}
\renewcommand{\author}{Dr.\,James Ashe}

\newcommand{\classNum}{137}

\newcommand{\classSec}{B}

\newcommand{\nameType}{}

\newcommand{\examType}{Final Review}

\renewcommand{\date}{\formatdate{16}{5}{2013}}

%%First page's header%%%%%%%%%%%%%%%%%%%%%%
\fancypagestyle{firststyle} %%header settings for first pagr
{
	\fancyhead[L]{MTH \classNum{}(\classSec{}) \author{}\\
	\textbf{\examType{}}}
	\fancyhead[C]{\textbf{The final will be given on \date~ 8-10am}}
	\setlength{\headsep}{14pt}
}
%%%%%%%%%%%%%%%%%%%%%%%%%%%%%%%%%%%%%%%%%%

%%Next page's header%%%%%%%%%%%%%%%%%%%%%%%
\setlist{nolistsep}
\lhead{\classNum{}(\classSec{})} 
\chead{\textbf{\nameType}} 
\cfoot{\thepage{}~of~\pageref{LastPage}} 
\setlength{\headsep}{5pt} 
%%%%%%%%%%%%%%%%%%%%%%%%%%%%%%%%%%%%%%%%%%%

%%Various settings; see the preamble for other settings%%%%%
%%\setenumerate[0]{leftmargin=12pt,itemindent=0pt,labelwidth=0pt}
\renewcommand{\labelenumi}{\textbf{\arabic{enumi}.}}
\renewcommand{\labelenumii}{\textbf{\alph{enumii}.}}
\thispagestyle{firststyle}
%%%%%%%%%%%%%%%%%%%%%%%%%%%%%%%%%%%%%%%%%%
\newcommand{\xmax}{10}
\newcommand{\xmin}{-10}
\newcommand{\xstep}{0} %must be xmin+n
\newcommand{\ymax}{10}
\newcommand{\ymin}{-10}
\newcommand{\ystep}{0} %must be ymin+n

\newcommand{\Dspace}{\vspace{.75cm}}\footnotesize

\subsubsection*{Solving equations}
\begin{itemize}
\item Solve an equation with fractions, ex.\textbf{ }$\frac{7x}{8}-2=3-\frac{15x}{16}$.

\begin{itemize}
\item Multipliy by the lowest common denominator and solve.
\end{itemize}
\item Solve an equation with absolute values, ex, $4|x+1|-7=6$.

\begin{itemize}
\item Isolate the absolute value part, e.g., $\left|x+1\right|$.
\item Solve for $\pm$ cases; Check your solutions.
\end{itemize}
\item Solve an equation with multiple unknown values for a given variable,
ex. $G=Q+(M-7)P$ for $M$.
\item Solve a $n^{\mbox{th }}$ root, ex. $\sqrt{x-1}+2=5$, $\sqrt[5]{x+1}=4$.

\begin{itemize}
\item Isoloate the radical and raise both sides to the power of $n$.
\item Solve and Check your answers.
\end{itemize}
\end{itemize}

\subsubsection*{Solving quadratics}
\begin{itemize}
\item Solve using the square root property, ex. \textbf{$(x-4)^{2}+1=2$.}

\begin{itemize}
\item Just solve like you do for $x$ and treat the other letters like numbers.
\item Isolate the squared part, e.g., $(x-4)^{2}$ and apply a square root
to both sides; don't forget the $\pm$.
\item Solve and check your answers.
\end{itemize}
\item Solve by factoring. ex. \textbf{$3x^{2}=-3x+6$.}

\begin{itemize}
\item Set the equation equal to zero and factor.
\item Find the value for $x$ that makes each factor $0$.
\end{itemize}
\item Solve using a graph. 

\begin{itemize}
\item Simply find the $x$ coordinates when the graph intersects (or touches)
the $x$-axis.
\item Note that this is an easy way to check the solutions for the above
problems using your calculator.
\end{itemize}
\end{itemize}

\subsubsection*{The circle: $\left(x-h\right)^{2}+\left(y-k\right)^{2}=r^{2}$ \textmd{where
$\left(h,k\right)$ is the center and $r$ is the radius.}}
\begin{itemize}
\item Find the center and radius of a circle from the standard equation,
ex. $(x-2)^{2}+(y+4)^{2}=16$.

\begin{itemize}
\item pick out your $h,k,$ and $r$; be careful with the $\pm$ signs.
\end{itemize}
\item Write the standard equation of a cirrcle for a given center and radius,
ex. $(9,4)$ and passing through the origin.

\begin{itemize}
\item Fill in the equation. If you're given a point it passes through instead
of the radius, you need to use the distance formula to find $r$.
\end{itemize}
\end{itemize}

\subsubsection*{Linear equations.}
\begin{itemize}
\item Find the slope of a line going through two points $(x_{1},y_{1})$
and $(x_{2},y_{2})$: $m=\frac{y_{2}-y_{1}}{x_{2}-y_{1}}$.
\item Find the slope of a line perpendicular to $\ell_{1}=m_{!}x+b_{1}$:
$m=-\frac{1}{m_{1}}$, or parallel to $\ell_{2}=m_{2}x+b_{2}$: $m=m_{2}$.
\item Find the equation for a line: $y-y_{1}=m(x-x_{1})+b$ where the line
passes through $(x_{1},y_{1})$.
\item Graph a line.
\end{itemize}

\subsubsection*{Complex numbers.}
\begin{itemize}
\item Perform $\times,\div,+,-$ with complex numbers (aka imaginary numbers),
ex., $(2-i)(4+8i)$.
\item Write a complex number in standard form, i.e., $a+bi$, ex., write
$\frac{4}{5+i}$ as $\frac{5}{6}-\frac{1}{6}i$.
\item Simply $i^{n}$, ex., $i^{3}=-i$, $i^{40}=1$,$i^{1}=i,\,i^{2}=-1,\,i^{3}=-i,\mbox{ and }i^{4}=1$.
Write $i^{n}$ in these terms, ex. $i^{42}=i^{40}\cdot i^{2}=\left(i^{4}\right)^{10}(-1)=1^{10}(-1)=-1$.
\end{itemize}

\subsubsection*{Functions.}
\begin{itemize}
\item Given a relation is a function, ex., $g=\left\{ \left(2,7\right),\left(30,-4\right),\left(45,7\right),\left(2,5\right)\right\} $.
Determine if it is a function, find its domain, and its range. 

\begin{itemize}
\item To determine if it is a function: Check that each $x$ , in $\left(x,y\right)$,
is paired with only one $y$ value. If any $x$ is paired with two
different $y's$ then the relation is NOT a function.
\item The domain is all the $x$ values, and the range is the $y$ values.
\end{itemize}
\item Given a graph, determine if it is a function, find its domain, and
its range. 

\begin{itemize}
\item Use the vertical line test, ex. see Exam2 solutions \#1.
\end{itemize}
\item Use transformations to graph a function.

\begin{itemize}
\item Know the parent funcitons of: $x^{2},\sqrt{x},a^{x},\log_{a}x$.
\item Know the transformations $f(x+a)\leftarrow$; $f(x-a)\rightarrow$;
$f(x)+a\uparrow$;$f(x)-a\downarrow$; $f(-x)$ $y$-axis reflection;
$-f(x)$ $x$-axis reflection. 
\end{itemize}
\item Difference quotient of a $f(x)$, ex., find the difference quotient
of $f(x)=7x+3$.

\begin{itemize}
\item Use the formula, $\frac{f(x+h)-f(x)}{h}$.
\item Find the average rate of change of a $f(x)$ from $a$ to $b$, ex.,
what is the average rate of change of $f(x)=x^{2}+1$ from $0$ to
$2$?
\end{itemize}
\item Find the inverse of a given function, ex., $f(x)=-\frac{1}{3}x+1$
find $f^{-1}$.
\item Compute function computation, ex., let \emph{$f(x)=x+4$ and $g(x)=x^{2}-x$
}find\emph{ $(f+g)(x)$ }and $(f\circ g)(2)$.
\end{itemize}

\subsubsection*{Quadratic functions.}
\begin{itemize}
\item Find the vertex, axis of symmetry, max/min, range, domain, directirx
and focus given a quadratic, ex., see \#1 and \#2 Exam2 and \#9 Exam3.

\begin{itemize}
\item First find the vertex, and then use it to answer the other parts.
\item If $a(x-h)^{2}+k$: the vertex is $(h.k)$. If $ax^{2}+bx+c$: the
vertex is $\left(\nicefrac{-b}{2a},f(\nicefrac{-b}{2a})\right)$.
\end{itemize}
\item Find a quadratic given the equation OR the directrix and focus. ex.,
see \#2 Exam3 and \#10 Exam4.

\begin{itemize}
\item Given the focus and directrix, $\frac{1}{4p}(x-h)^{2}+k$ where $p$
is $\nicefrac{1}{2}$ the distance between the directrix and the focus.
If the focus is $(x,y)$ then the vertex is $(x,y-p)$.
\end{itemize}
\end{itemize}

\subsubsection*{Polynomials.}
\begin{itemize}
\item Use traditional or synthetic division, ex., let $P(x)=x^{2}-5x+2$

\begin{itemize}
\item find the quotient and remainder, ex., find the quotient $Q(x)$ and
remainder $r$ when $P(x)$ is divided by $x-2$.
\item Determine if $a$ is a root of $P(x)$; $a$ is a root if and only
if $P(a)=0$ or $P(x)\div(x-a)$ has no remainder.
\item Write $\frac{P(x)}{x-a}$, as a sum of polynomial and rational function:
use the formula $\frac{P(x)}{x-a}=Q(x)+\frac{r}{x-a}$.
\end{itemize}
\item Find the simplest polynomial with the given roots, ex., for roots
$-2,2,5$ it is $(x+2)(x-2)(x-5)$.
\item Find the $x$-ints of a polynomial; set the function equal to $0$
and solve.
\item Determine if the $x$-ints pass or touches the $x$-axis.

\begin{itemize}
\item look at the multiplicity: if it is even it touches, if it is odd it
passes.
\end{itemize}
\item Find all vertical and horizontal asymptotes.

\begin{itemize}
\item Vertical: where the function is undefined (usually when it is $0$
in the denominator), ex., $\frac{2x+6}{x-2}$ has vertical asymptote
at $x=2$ since this makes the denominator $0$.
\item Horizontal: exists if the limit of the function converges, ex., $\lim_{x\rightarrow\infty}\frac{2x+6}{x-2}\cdot\frac{\nicefrac{1}{x}}{\nicefrac{1}{x}}=\frac{2+\nicefrac{6}{x}}{1-\nicefrac{2}{x}}=2$
so $\frac{2x+6}{x-2}$ has horizontal asymptote $2$.
\end{itemize}
\end{itemize}

\subsubsection*{Limits.}
\begin{itemize}
\item Determine $y$ as $x\rightarrow\pm\infty$, ex., $x^{2}-20$ as $x\rightarrow-\infty$
\item Determine ${\displaystyle \lim_{x\rightarrow a^{\pm}}}f(x)$, for
polynomial, rational, exponential, and lograthmic functions, ex.,
${\displaystyle \lim_{x\rightarrow\infty}\frac{1}{x^{3}}}$ and ${\displaystyle \lim_{x\rightarrow2^{-}}\frac{3}{x-2}}$

\begin{itemize}
\item ${\displaystyle \lim_{x\rightarrow\infty}x=\infty}$ and ${\displaystyle \lim_{x\rightarrow\infty}\frac{1}{x}=0}$,
${\displaystyle \lim_{x\rightarrow0^{-}}\frac{1}{x}=-\infty}$, and
${\displaystyle \lim_{x\rightarrow0^{+}}\frac{1}{x}=\infty}$
\end{itemize}
\end{itemize}

\subsubsection*{Exponential and Logaritthmic functions.}
\begin{itemize}
\item Evaluate exponents, ex., $-2^{-3}$, $16^{\nicefrac{2}{4}}$, $\ln(e^{5})$,
$\log_{b}1$.

\begin{itemize}
\item Need to be familiar with the rules of exponents and logs.
\end{itemize}
\item Find the graph, domain, and range given a log or exponential function.
\item Use properties of logs to condense an expression, ex., $\ln2+\ln x^{6}=\ln\left(2x^{6}\right)$.
\item Use the properties of logs to expand an expression, ex., $\log\left(\frac{10^{6}}{x}\right)=6-\log x$.
\item Solve an exponential function, ex., $1500=500e^{.0.10\cdot t}$ for
$t$.

\begin{itemize}
\item Use a $\log_{b}x$, usually $\ln x$, on BOTH sides ex., $\ln1500=\ln500e^{.0.10\cdot t}$
\end{itemize}
\item Solve a logarithmic function, ex., $\ln(40)=3-\ln(x)$ for $x$.

\begin{itemize}
\item Use an exponent of the same base as the logs on BOTH sides, ex., $\log_{6}3x=2\Rightarrow6^{\log_{6}3x}=6^{2}\Rightarrow3x=36$.
\end{itemize}
\end{itemize}

\end{document}
